% Replace the entire Discussion body, including text outside old generated markers.
% First complete-results scientific writing pass; independent review still required.

The completed comparison supports conditional predictive benefits rather
than general mixed-forest superiority. Several shallow fixed replacement
contrasts retain small joint-Holm signed-rank values, including both
all-feature replacements and the square-root-feature horizon-three
replacement. Their positive means coexist with zero or small medians
and substantial extra fitting work. Uniformly increasing horizon is
especially unfavorable for the unrestricted all-feature models, whereas
square-root-feature forests show a different pattern. Better local
terminal training impurity is not a certificate of better prediction,
and neither this depth pattern nor a difference between feature policies
identifies a causal mechanism.

Training-only selection changes the practical comparison. The inclusive
family has a positive estimated mean advantage over selected RF, with a
pointwise mean-effect interval above zero. Such a mean gain may be practically
relevant, but its adjusted signed-rank test
does not reject the location null, its median effect is zero, and the
family-weighted intervals and those for the 38 tasks outside the named
synthetic/constructed set cross zero. These summaries answer different questions. They permit a heterogeneous mean
benefit on the full cohort, not a claim of no improvement, equivalence,
or reliable improvement for a typical future task. Because the expanded
families often select pure forests, their aggregate effects evaluate
larger training-selected candidate spaces, not the isolated benefit of
sparse mixtures. Both mixed-capable spaces contain the entire RF space,
but this nesting does not guarantee a held-out improvement. Their larger
selection spaces and search work are part of the comparison.

The shallow experiment remains a capacity-constrained question, not a
recommended universal forest depth. Tree-size regularization studies
motivate conditional benefits of smaller members, particularly in noisy
regression settings \cite{zhou2023trees,duroux2018impact}; they do not
identify depth three as optimal for these bootstrap Gini classifiers.
The fixed CART forests have higher mean accuracy at deeper capacities,
and selected RF chooses depth six or unrestricted depth in most outer
tasks. Thus deeper forests remain relevant and often useful controls.
Growing-tree consistency results do not establish consistency at fixed
depth three \cite{scornet2015consistency}, and shallow boosting results
do not justify an independently fitted bootstrap forest.

Computation remains a substantial limitation. Full standalone selection
work is much larger for the expanded families than for RF, particularly
when horizon-three banks are required, even if the selected refit is
relatively cheap. The observed costs characterize the implementation,
candidate space, and shared hardware load, not every non-greedy method.
Local lookahead is established \cite{murthy1995lookahead,donick2021stepwise};
the additive child recursion gives a worst-case candidate-evaluation
bound exponential in local horizon, rather than the number of complete
split assignments. Final fitting also sums searches over committed
nodes. Fixed horizon is distinct from final depth. Prefix-count reuse
and Cython improve execution without introducing a new optimization
principle or a global certificate. Cached, bounded, and hybrid solvers,
including SPLIT, require objective- and constraint-aligned comparisons
before an accuracy or efficiency ranking is justified
\cite{aglin2020caching,demirovic2022murtree,vanderlinden2023separable,babbar2025nearoptimal}.

The sample is small, previously explored, and includes related and
constructed tasks. Repeated overlapping holdouts are not independent
datasets; known-family weighting addresses only documented dependence.
Complete-case deletion, supplied numeric encodings, rare-class validation
fallbacks, and bounded tuning grids limit the conclusions. Matching common
hyperparameters does not create equal computational budgets, and remaining
threshold, tie, and random-feature-order differences preclude bitwise
algorithmic identity. The descriptive mean-cost curves neither enforce a
per-dataset budget nor account for selection work. No causal diversity
effect, validated resource frontier, large-data scalability, or production
readiness has been established. A prospective resource-constrained study
would need training-only budget selection, final-fit feasibility checks,
complete cost accounting, and larger independent tasks.
